%---------------------------------------------------------------%
\addchap{AGRADECIMENTOS}
\begin{refsection}

Gostaríamos de expressar nossos mais sinceros agradecimentos a todas as pessoas
que contribuíram para a produção deste livro. Em especial, agradecemos ao
Excelentíssimo Governador do Estado do Pará, \textbf{Helder Zaluth Barbalho},
por sua visão estratégica voltada ao desenvolvimento sustentável da região amazônica,
aliando crescimento econômico, inclusão social e preservação ambiental. \vskip0.3cm

Seu compromisso com a valorização da Ciência e com o fortalecimento das políticas
públicas baseadas em evidências reforça o papel fundamental que a Estatística exerce
na construção de um futuro mais justo, eficiente e sustentável para o Estado do Pará.
\vskip0.3cm

Manifestamos também nossa gratidão à \textbf{Secretaria de Segurança Pública do Estado do Pará},
pelo apoio institucional concedido a este projeto. Seu compromisso com a sociedade e
com o aprimoramento contínuo da gestão pública é inestimável. Esperamos que este
compêndio possa servir como um recurso valioso para servidores e profissionais que
buscam capacitação na tomada de decisões fundamentadas em dados.
\vskip0.3cm

Registramos ainda nosso agradecimento à \textbf{Autarquia de Trânsito do Estado do Pará (DETRAN-PA)},
que, sob a gestão de sua Diretora-Geral, Dra. \textbf{Renata Mirella Freitas Guimarães de Sousa Coelho},
gentilmente possibilitou o uso de suas dependências para importantes discussões técnicas
e etapas essenciais da produção desta obra.
\vskip0.3cm

Agradecemos de forma especial aos Estatísticos
\textbf{Alexandre Rodrigues Ramos}, Analista Judiciário do Tribunal de Justiça do Estado do Pará;
\textbf{Dennison Célio de Oliveira Carvalho}, Professor Doutor Associado da Universidade Federal
do Oeste do Pará; e
\textbf{Paulo Cerqueira dos Santos Junior}, Professor Adjunto da Universidade Federal do Pará,
pelo generoso compartilhamento de seus conhecimentos e pela revisão criteriosa deste material.

\newpage
Suas contribuições, sugestões e correções foram essenciais para assegurar a qualidade,
a precisão técnica e a clareza pedagógica do conteúdo apresentado.
\vskip0.3cm

Sem a dedicação e o conhecimento compartilhados por todos esses profissionais,
esta obra não teria sido possível. Seus esforços enriquecem cada capítulo deste livro,
proporcionando aos leitores uma compreensão aprofundada das ferramentas estatísticas
necessárias para a formulação de decisões informadas no setor público.
\vskip0.3cm

Agradecemos igualmente aos nossos atuais e ex-alunos de graduação e pós-graduação,
cuja colaboração e entusiasmo contribuíram em diferentes momentos para a construção
deste trabalho. Estendemos nossa gratidão aos amigos e colegas, pelo incentivo,
pela leitura atenta e pelos valiosos comentários ao longo do processo de escrita.
\vskip0.3cm

Este livro é dedicado a todas as pessoas que, assim como nós, são apaixonadas por
Estatística, Matemática e pela linguagem de programação \textit{R}, e que acreditam
na importância de levar o conhecimento científico às novas gerações.

\printbibliography[heading=subbibliography]

\end{refsection}

%---------------------------------------------------------------%
\addchap{EPÍGRAFE}

\vspace{5cm}

\begin{minipage}{10cm}

\textit{``Se a Matemática é o pincel com que Deus desenhou o universo,
então a Estatística é a ferramenta humana criada para tentar compreendê-lo.''}

\vspace{0.3cm}

\hfill (\textbf{Waldenei Travassos Queiroz, 2012})

\vspace{5cm}

\hfill \textbf{Professor Titular Aposentado da UFRA}

\end{minipage}
%---------------------------------------------------------------%
